% =================================================
% Содержание темы
% Этот файл подключается и к цветной, и к чёрно-белой версии.
% Здесь не должно быть настроек оформления.
% =================================================

\begin{center}
			{\fontsize{25}{30}\selectfont\bfseries\sffamily\color{Primary}
				Геометрическое место точек}
		\end{center}
	
	\vspace{0.5em}
	
	\section{Что такое геометрическое место точек}
	
		Многие геометрические фигуры можно рассматривать как \textbf{множества точек}.
		Например, прямую, окружность и другие геометрические фигуры можно
		рассматривать как множества точек.
	
		Часто точки некоторой фигуры обладают общим свойством. Например, все точки
		окружности равноудалены от её центра.
	
	Пусть дана точка \(O\). Будем отмечать на плоскости все точки, которые
	находятся от неё на расстоянии \(3\) см.
	
	Если отметить только несколько таких точек, фигура ещё не получится.
	Но если отметить все точки, удовлетворяющие условию, они образуют окружность
	с центром \(O\) и радиусом \(3\) см.

	\gmtcirclefigure
	
	При этом важно не пропустить ни одной подходящей точки и не добавить точек,
	находящихся от \(O\) на другом расстоянии.
	
	Именно поэтому геометрическое место должно удовлетворять двум условиям:
	
	\begin{enumerate}
		\item каждая точка найденной фигуры обладает заданным свойством;
		\item любая точка, обладающая заданным свойством, принадлежит найденной фигуре.
	\end{enumerate}
	
	\begin{definitionbox}
		\textbf{Геометрическим местом точек}, или сокращённо \textbf{ГМТ},
			называется множество всех точек, которые обладают заданным свойством.
	\end{definitionbox}
	
	\Needspace{8\baselineskip}
	\section{Точки на заданном расстоянии от данной точки}
	
	Пусть на плоскости дана точка \(O\) и положительное число \(R\).
	Рассмотрим все точки \(M\), для которых
	
	\[
	OM=R.
	\]
	
	Это означает, что расстояние от точки \(M\) до точки \(O\) равно \(R\).
	
	Все такие точки образуют окружность с центром \(O\) и радиусом \(R\).
	
	\begin{propertybox}
		Геометрическим местом точек, находящихся на расстоянии \(R\)
		от данной точки \(O\), является окружность с центром \(O\)
		и радиусом \(R\).
	\end{propertybox}
	
	Почему это действительно геометрическое место?
	
	\begin{itemize}
		\item Каждая точка окружности находится на расстоянии \(R\) от её центра.
		\item Любая точка, находящаяся на расстоянии \(R\) от точки \(O\),
		лежит на этой окружности.
	\end{itemize}
	
	\subsection{Как изменение условия меняет фигуру}
	
	Если изменить условие, изменится и геометрическое место точек.

	\gmtdistancecomparisonfigure
	
	\begin{itemize}
			\item Условие \(OM<R\) задаёт внутренность круга без его границы.
		\item Условие \(OM\leq R\) задаёт круг, то есть все точки внутри окружности
		вместе с точками самой окружности.
		\item Условие \(OM>R\) задаёт все точки, лежащие вне круга.
	\end{itemize}
	
	\begin{attentionbox}[Обратите внимание]
		Знаки \(<\), \(=\) и \(\leq\) задают разные множества точек.
	\end{attentionbox}
	
	\Needspace{8\baselineskip}
	\section{Точки, равноудалённые от концов отрезка}
	
	Пусть даны две различные точки \(A\) и \(B\).
	
	Рассмотрим все точки \(M\), для которых
	
	\[
	MA=MB.
	\]
	
		Такие точки равноудалены от точек \(A\) и \(B\).
	
	Геометрическим местом этих точек является прямая, проходящая через середину
	отрезка \(AB\) перпендикулярно ему.
	
	Эта прямая называется \textbf{серединным перпендикуляром к отрезку \(AB\)}.

	\gmtperpendicularbisectorfigure

	\begin{propertybox}
		Геометрическим местом точек, равноудалённых от концов отрезка,
		является серединный перпендикуляр к этому отрезку.
	\end{propertybox}
	
	Если точка \(M\) лежит на серединном перпендикуляре к отрезку \(AB\), то
	
	\[
	MA=MB.
	\]
	
	Верно и обратное: если для некоторой точки \(M\) выполняется равенство
	
	\[
	MA=MB,
	\]
	
	то точка \(M\) лежит на серединном перпендикуляре к отрезку \(AB\).
	
	Середина самого отрезка \(AB\) также принадлежит этому геометрическому месту.
	
	\Needspace{8\baselineskip}
	\section{Точки внутри угла, равноудалённые от его сторон}
	
	\begin{notebox}[Напомним]
			Расстоянием от точки до прямой называется длина отрезка перпендикуляра,
			проведённого из данной точки к прямой.
	\end{notebox}

	\gmtdistancetolinefigure

		Пусть дан угол. Рассмотрим все точки \(M\), лежащие внутри этого угла
		и равноудалённые от его сторон.
		
		Геометрическим местом таких точек является \textbf{биссектриса данного угла}.

	\gmtbisectorfigure

	\begin{propertybox}
		Геометрическим местом точек, лежащих внутри угла и равноудалённых
		от его сторон, является биссектриса этого угла.
	\end{propertybox}
	
	\begin{notebox}
		Важно, что в данном утверждении рассматриваются только точки,
		расположенные внутри угла.
	\end{notebox}
	
		Каждая точка биссектрисы равноудалена от сторон угла.
	И наоборот, любая точка внутри угла, равноудалённая от его сторон,
	лежит на биссектрисе.
	
	\Needspace{8\baselineskip}
	\section{Почему в ГМТ не должно быть лишних и пропущенных точек}
	
	Рассмотрим условие
	
	\[
	OM=R.
	\]
	
	Ему удовлетворяют все точки окружности с центром \(O\) и радиусом \(R\).
	
	Можно ли считать искомым геометрическим местом только полуокружность?
	
	Нет. Каждая точка полуокружности действительно находится на расстоянии \(R\)
		от точки \(O\). Но на другой полуокружности также лежат точки,
	для которых выполняется условие \(OM=R\). Эти точки оказались пропущены.
	
	Можно ли считать искомым геометрическим местом весь круг?
	
	Тоже нет. Внутри круга находятся точки, для которых
	
	\[
	OM<R.
	\]
	
	Они не удовлетворяют условию \(OM=R\), поэтому являются лишними.

	\gmtextraandmissingpointsfigure
	
	\begin{importantbox}[Лишние и пропущенные точки]
		Полуокружность содержит не все подходящие точки, а круг содержит лишние точки.
		Искомым геометрическим местом является именно окружность:
		на ней нет лишних точек и ни одна подходящая точка не пропущена.
	\end{importantbox}
	
	\Needspace{8\baselineskip}
	\section{Как доказать, что фигура является геометрическим местом точек}
	
	Чтобы доказать, что найденная фигура является ГМТ, нужно проверить два условия.
	
		\subsection{Прямое утверждение: у найденной фигуры нет лишних точек}
	
	Нужно взять любую точку \(M\) этой фигуры и доказать,
	что она обладает заданным свойством.
	
	Например, если предполагаемым ГМТ является серединный перпендикуляр
	к отрезку \(AB\), нужно взять любую его точку \(M\) и доказать, что
	
	\[
	MA=MB.
	\]
	
		\subsection{Обратное утверждение: ни одна подходящая точка не пропущена}
	
	Нужно взять любую точку \(M\), обладающую заданным свойством,
	и доказать, что она принадлежит найденной фигуре.
	
	В примере с серединным перпендикуляром нужно доказать:
	если
	
	\[
	MA=MB,
	\]
	
	то точка \(M\) лежит на серединном перпендикуляре к отрезку \(AB\).
	
	\begin{notebox}[Итог]
		Только после проверки обоих условий можно утверждать,
		что найденная фигура является геометрическим местом точек.
	\end{notebox}
	
	\Needspace{8\baselineskip}
	\section{Несколько условий и пересечение геометрических мест}
	
	Одна и та же точка может одновременно удовлетворять нескольким условиям.
	
	В таком случае каждому условию соответствует своё геометрическое место,
		а искомые точки принадлежат пересечению этих геометрических мест.
	
	\begin{examplebox}
		Найдём все точки \(M\), для которых одновременно выполняются условия
		
		\[
		MA=4\text{ см}
		\qquad\text{и}\qquad
		MB=3\text{ см}.
		\]
		
		Первому условию соответствует окружность с центром \(A\)
		и радиусом \(4\) см.
		
		Второму условию соответствует окружность с центром \(B\)
		и радиусом \(3\) см.
		
		Искомые точки должны принадлежать обеим окружностям.
		Следовательно, они являются точками пересечения этих окружностей.
	\end{examplebox}

	\gmtcirclesintersectionfigure

		В зависимости от расположения точек \(A\) и \(B\) окружности могут:
		
		\begin{itemize}
			\item иметь две общие точки;
			\item иметь одну общую точку;
			\item не иметь общих точек.
		\end{itemize}
		
		Таким образом, геометрические места точек позволяют переводить словесные
		условия на язык известных геометрических фигур.
	
	\Needspace{8\baselineskip}
	\section{Каким может быть геометрическое место точек}
	
	Геометрическое место точек не обязательно является прямой или окружностью.
	
	Оно может представлять собой:
	
	\begin{itemize}
		\item одну точку;
		\item несколько отдельных точек;
		\item прямую или луч;
		\item окружность;
		\item некоторую область;
		\item пустое множество.
	\end{itemize}
	
	Например, точек \(M\), для которых одновременно выполняются условия
	
	\[
	MA=2\text{ см}
	\qquad\text{и}\qquad
	MA=3\text{ см},
	\]
	
	не существует.
	
	В этом случае геометрическое место точек является пустым множеством.
	
	\Needspace{10\baselineskip}
	\section{Основные геометрические места точек}
	
	\renewcommand{\arraystretch}{1.5}
	
	\begin{center}
		\begin{tabularx}{\textwidth}{
				>{\raggedright\arraybackslash}p{0.29\textwidth}
				>{\raggedright\arraybackslash}X
			}
			\toprule
			\textbf{Свойство точки \(M\)}
			&
			\textbf{Геометрическое место точек}
			\\
			\midrule
			
			\(OM=R\)
			&
			Окружность с центром \(O\) и радиусом \(R\)
			\\
			
			\(OM<R\)
			&
			Внутренность круга без границы
			\\
			
			\(OM\leq R\)
			&
			Круг с центром \(O\) и радиусом \(R\)
			\\
			
			\(OM>R\)
			&
			Множество точек вне круга
			\\
			
			\(MA=MB\)
			&
			Серединный перпендикуляр к отрезку \(AB\)
			\\
			
			Точка \(M\) лежит внутри угла и равноудалена от его сторон
			&
			Биссектриса угла
			\\
			
			\bottomrule
		\end{tabularx}
	\end{center}
	
	\Needspace{10\baselineskip}
	\section{Как решать задачи на геометрическое место точек}
	
	\begin{algorithmbox}
		\begin{enumerate}
			\item Обозначить переменную точку буквой \(M\).
			\item Записать свойство, которому должна удовлетворять точка \(M\).
			\item Вспомнить, какая известная фигура задаётся этим свойством.
			\item Проверить, нет ли в найденной фигуре лишних точек.
			\item Убедиться, что все точки с нужным свойством найдены.
				\item Если задано несколько условий, построить соответствующие им
				геометрические места и найти их пересечение.
		\end{enumerate}
	\end{algorithmbox}
	
	\begin{propertybox}[Главная идея]
		\[
		\text{свойство точки}
		\quad\longleftrightarrow\quad
		\text{геометрическая фигура}.
		\]
	\end{propertybox}
	
	Задачи на геометрические места точек учат переводить условие задачи
	на язык геометрических фигур и находить точки, удовлетворяющие одному
	или нескольким требованиям.
	
	\begin{questionsbox}
		\begin{enumerate}
			\item Что называется геометрическим местом точек?
			\item Какие два утверждения необходимо доказать, чтобы подтвердить,
			что найденная фигура является ГМТ?
			\item Какое геометрическое место задаётся условием \(OM=5\text{ см}\)?
			\item Какое геометрическое место задаётся условием \(OM\leq 5\text{ см}\)?
			\item Где находятся точки, равноудалённые от концов отрезка?
			\item Где находятся точки внутри угла, равноудалённые от его сторон?
			\item Почему полуокружность не является ГМТ точек,
			находящихся на расстоянии \(R\) от точки \(O\)?
			\item Как находят точки, которые должны удовлетворять сразу двум
			геометрическим условиям?
		\end{enumerate}
	\end{questionsbox}
